\bookchapter{Capítulo 24: Análise aprofundada de permissões Root e bloqueio do Bootloader}

\chapterintro{Este capítulo é uma extensão aprofundada do conteúdo do Capítulo 20, adequado para usuários que desejam realmente entender os mecanismos de baixo nível do celular e até mesmo tentar flashing de ROM. Atenção: essas operações envolvem certos riscos, proceda com cautela.}

\booksection{I. Compreensão aprofundada do Bootloader}

O Bootloader é o primeiro programa executado quando o celular é ligado. Sua função é inicializar o hardware e carregar o kernel do sistema operacional. Pode-se entender o Bootloader como o “guardião” do celular.

\booksubsection{Como funciona o bloqueio do Bootloader:}

\begin{itemize}
  \item \textbf{Verificação de assinatura}: o bloqueio do Bootloader verifica a assinatura digital da imagem de inicialização. Apenas sistemas assinados pelo fabricante podem iniciar.
  \item \textbf{Prevenção de downgrade}: impede que o usuário faça flash de versões antigas do sistema para explorar vulnerabilidades conhecidas.
  \item \textbf{Proteção de dados}: mesmo com um novo sistema instalado, sem a senha da conta original, os dados do usuário permanecem inacessíveis.
\end{itemize}

\booksubsection{Armazenamento e verificação em nível de hardware}

O estado de bloqueio do Bootloader (Bloqueado/Desbloqueado) e a chave pública do fabricante geralmente são armazenados em áreas de hardware protegidas (como a partição devinfo da Qualcomm, eFuse ou partição RPMB), para evitar adulteração por software. A cada inicialização, o Bootloader lê o estado desta área segura e, com base nisso, decide se executa a verificação rigorosa de assinatura.

\booksubsection{Verificação de assinatura digital}

Ao publicar atualizações de sistema ou imagens, o fabricante as assina com sua chave privada. Na inicialização do dispositivo, o Bootloader usa a chave pública incorporada para verificar as imagens do sistema (como boot.img, recovery.img). O processo funciona assim:

\begin{enumerate}
  \item Recalcular o valor de hash da imagem do sistema.
  \item Descriptografar a assinatura incluída na imagem usando a chave pública, obtendo um valor de hash.
  \item Comparar os dois valores de hash; se forem iguais, a verificação é aprovada. Isso garante que a imagem do sistema não foi adulterada por terceiros.
\end{enumerate}

\booksubsection{Cadeia de inicialização segura}

O bloqueio do Bootloader é o primeiro elo de toda a cadeia de inicialização segura. Somente após sua verificação ser aprovada, o próximo estágio do bootloader (como o kernel) é carregado e executado, garantindo assim, camada por camada, a confiabilidade do ambiente de inicialização do sistema.

\booksubsection{Bloqueado vs. Desbloqueado: o impacto duplo do bloqueio do Bootloader}

Para entender mais intuitivamente os prós e contras do bloqueio do Bootloader, apresentamos uma tabela comparativa em 5 dimensões principais:

\begin{center}
\begin{tabularx}{\linewidth}{|X|X|X|}
\hline
\textbf{Função principal} & \textbf{Estado bloqueado (padrão)} & \textbf{Estado após desbloqueio} \tabularnewline
\hline
Segurança do sistema & Impede a instalação de sistemas não oficiais não assinados, protegendo a integridade e a segurança de baixo nível do sistema\newline Alta segurança: carrega apenas imagens com assinatura oficial, resistindo eficazmente à adulteração por firmware malicioso & Segurança reduzida: permite instalação de ROMs de terceiros, o mecanismo de verificação de integridade do sistema fica inativo; malware com privilégios Root pode adulterar profundamente o sistema \tabularnewline
Estabilidade do sistema & Garante que o celular execute o sistema oficial testado e otimizado pelo fabricante, assegurando estabilidade e fluidez\newline Estável: usa firmware oficial, funcionalidades completas, atualizações OTA confiáveis & Risco aumentado: ROMs de terceiros podem ter problemas de compatibilidade, causando travamentos, brick, mau funcionamento de hardware (câmera, carregamento rápido), e falhas em atualizações OTA \tabularnewline
Direitos de garantia & O fabricante fornece sistema oficial estável com o Bootloader bloqueado para oferecer serviço de garantia\newline Garantia válida: desfrute do serviço de garantia oficial gratuito (danos não causados por mau uso) & Garantia anulada: a grande maioria dos fabricantes considera o desbloqueio do Bootloader como renúncia voluntária à garantia; mesmo ao rebloquear, o histórico de desbloqueio pode ser detectado e a garantia recusada \tabularnewline
Segurança de dados & O processo de desbloqueio geralmente força a exclusão de dados, impedindo que dados antigos sejam explorados maliciosamente no novo ambiente\newline Dados seguros: no uso diário, o bloqueio do Bootloader limita a capacidade de programas maliciosos modificarem a partição do sistema & Risco de dados: o desbloqueio força a exclusão de todos os dados do usuário; se não criptografado corretamente após o desbloqueio, o armazenamento pode ficar inseguro; em caso de perda, é mais fácil extrair dados por flashing \tabularnewline
Restrições funcionais & Garante o funcionamento normal de funções que dependem de ambiente seguro (como pagamentos, apps bancários)\newline Funcionalidade completa: todas as funções oficiais operam normalmente & Funcionalidade limitada: alguns apps sensíveis à segurança (como Google Pay, apps bancários) podem se recusar a executar ao detectar o estado de desbloqueio do Bootloader; serviços DRM podem ser rebaixados, afetando a reprodução de vídeo em alta definição \tabularnewline
\hline
\end{tabularx}
\end{center}

\booksection{II. Processo detalhado de desbloqueio do Bootloader (exemplo Xiaomi)}

\begin{enumerate}
  \item \textbf{Solicitar permissão de desbloqueio}: solicite o desbloqueio no site da ferramenta de desbloqueio Xiaomi.
  \item \textbf{Vincular conta}: vincule a conta Xiaomi nas configurações do celular.
  \item \textbf{Período de espera}: dependendo do nível da conta, é necessário aguardar de 7 a 30 dias.
  \item \textbf{Operação de desbloqueio}: use a ferramenta de desbloqueio Xiaomi conectada ao computador para desbloquear.
\end{enumerate}

\booksection{III. Diversas soluções de Root}

\begin{itemize}
  \item \textbf{Magisk}: solução moderna de Root, design Systemless, capaz de ocultar o estado de Root.
  \item \textbf{Kernelsu}: solução de Root baseada no kernel, mais baixo nível, melhor ocultação.
  \item \textbf{SuperSU}: solução de Root tradicional, atualmente descontinuada.
\end{itemize}

\booksection{IV. Introdução ao flashing de ROM}

Flashing de ROM refere-se a instalar um novo sistema operacional no celular. Tipos comuns de flashing:

\begin{itemize}
  \item \textbf{Flash de ROM oficial}: instalar o sistema oficial, usado para recuperar brick ou downgrade.
  \item \textbf{Flash de Recovery de terceiros}: como TWRP, usado para backup/restauração e instalação de ROMs de terceiros.
  \item \textbf{Flash de ROM de terceiros}: como LineageOS, Pixel Experience, etc., para experimentar sistemas diferentes.
\end{itemize}

\booksection{V. Samsung Knox: o custo irreversível do desbloqueio do Bootloader}

Diferente de outros fabricantes Android, o mecanismo de desbloqueio do Bootloader da Samsung possui um design exclusivo de “fusível”:\textbf{para celulares Samsung, o desbloqueio aciona o fusível Knox, perdendo permanentemente algumas funções de segurança (como Samsung Pay); apaga todos os dados protegidos; afeta a garantia}. Especificamente:

\begin{itemize}
  \item \textbf{O fusível é de uso único}: o celular possui um eFuse interno. Uma vez detectado o desbloqueio do Bootloader, ele queima permanentemente, sem possibilidade de restauração.
  \item \textbf{Funcionalidades afetadas}: Samsung Pay (cartão de transporte/cartão bancário), pasta segura Knox, gerenciamento MDM empresarial, e alguns apps de saúde e seguros deixarão de funcionar.
  \item \textbf{Limpeza forçada de dados}: após o fusível ser acionado, o celular força a restauração de fábrica, apagando todos os dados locais (dados na nuvem não são afetados).
  \item \textbf{Impacto na garantia}: a Samsung declara oficialmente que, após o disparo do Knox, o serviço de garantia não é mais aplicável.
\end{itemize}

\begin{tcbwarn}{Pense duas vezes antes de agir}
O fusível Knox da Samsung é\textbf{irreversível}. Uma vez acionado, não há como voltar ao estado original de fábrica. Se você deseja apenas usar funcionalidades que originalmente exigiriam Root, verifique primeiro se existem alternativas que não necessitem de Root.

\end{tcbwarn}

\textbf{Aviso importante}\par

Flashing de ROM envolve riscos, proceda com cautela. Pode resultar em:

\begin{itemize}
  \item Celular brickado (não liga)
  \item Perda da garantia
  \item Perda de dados
  \item Mau funcionamento de funções de pagamento
\end{itemize}

Antes de fazer o flashing, faça um backup completo, leia atentamente os tutoriais e confirme se o modelo e a versão são compatíveis.

\begin{tcbinfo}{Recursos recomendados}
\begin{itemize}
  \item \textbf{XDA Developers}: a maior comunidade de desenvolvedores Android do mundo, com abundantes tutoriais de flashing e recursos de ROM.
  \item \textbf{CoolAPK}: comunidade digital chinesa renomada, com muitos tutoriais e discussões em chinês.
  \item \textbf{Fóruns específicos por modelo}: como comunidade Xiaomi, comunidade OnePlus, etc., com tutoriais oficiais e da comunidade.
\end{itemize}

\end{tcbinfo}

\booksection{VI. Perguntas frequentes (Q\& A)}

\textbf{Q1: Após desbloquear o Bootloader, é possível restaurá-lo?}\par

\textbf{R:}Sim, é possível rebloquear. Geralmente é necessário fazer o flash do pacote de recuperação oficial completo (fastboot) e usar um script de flashing específico (como\texttt{flash\_all\_lock.\allowbreak bat}da Xiaomi) para rebloquear. Mas atenção: após rebloquear, alguns fusíveis de segurança acionados pelo desbloqueio (como o Knox da Samsung) podem ser irreversíveis.

\textbf{Q2: Como obter Root após o desbloqueio?}\par

\textbf{R:}Desbloquear o Bootloader é apenas o primeiro passo. O Root geralmente requer:

\begin{enumerate}
  \item Instalar um Recovery personalizado (como TWRP).
  \item Através do Recovery, instalar um gerenciador de Root como Magisk para corrigir a imagem de boot e obter privilégios Root.
\end{enumerate}
